% ---------------------------------------------------------------------
%  CHƯƠNG 2 — CƠ SỞ LÝ THUYẾT VÀ CÔNG NGHỆ
% ---------------------------------------------------------------------
\chapter{Cơ sở lý thuyết và công nghệ}
\label{ch:coso}

\begin{muctieu}
Chương này trình bày nền tảng của bài toán nhận diện quần áo, gợi ý phối đồ,
thử trang phục ảo và các nhóm công nghệ được xác định trong phạm vi đồ án.
\end{muctieu}

\section{Kiến thức nền tảng}

\begin{dinhnghia}[Phát hiện trang phục]
Phát hiện trang phục là bài toán xác định hộp giới hạn, nhãn lớp và độ tin
cậy của một hoặc nhiều món đồ xuất hiện trong ảnh. Với mỗi dự đoán, mô hình
trả về hộp $B_p$, nhãn $y \in \mathcal{Y}$ và một điểm tin cậy; trong đồ án,
$\mathcal{Y}$ gồm 21 lớp quần áo và phụ kiện.
\end{dinhnghia}

YOLO thực hiện phát hiện đối tượng trong một lần suy luận, kết hợp định vị và
nhận dạng lớp trong cùng mạng nơ-ron~\cite{Redmon2016}. Notebook của đồ án
khởi tạo từ checkpoint \code{yolo26n.pt}~\cite{UltralyticsYOLO26}, sau đó fine-tune tuần tự qua các
phiên bản V1, V2 và V3 bằng thư viện Ultralytics. Ảnh đầu vào được đưa về
kích thước $640\times640$ pixel; cấu hình dữ liệu giữ nguyên thứ tự 21 lớp để
checkpoint và API diễn giải nhãn nhất quán.

Mức chồng lấp giữa hộp dự đoán $B_p$ và hộp thật $B_g$ được đo bằng IoU ở
Công thức~\eqref{eq:iou-report}:
\begin{equation}
\mathrm{IoU}(B_p,B_g)=\frac{|B_p\cap B_g|}{|B_p\cup B_g|}.
\label{eq:iou-report}
\end{equation}
IoU được dùng khi ghép dự đoán với nhãn thật và khi tính precision, recall,
mAP@0,5 và mAP@0,5:0,95. Các đại lượng này phù hợp với bài toán phát hiện hơn
accuracy của bài toán phân loại ảnh~\cite{COCO2014}.

DeepFashion In-shop cung cấp 52.712 ảnh cùng hộp giới hạn và thông tin sản
phẩm~\cite{DeepFashion2016}. Notebook V3 chuyển hộp giới hạn sang định dạng
YOLO, chia dữ liệu theo \code{item\_id} với hạt giống 42 để các góc chụp của
cùng một sản phẩm không rơi vào nhiều tập, rồi trộn với bộ dữ liệu Roboflow
ban đầu. Phép tăng cường trong cấu hình huấn luyện gồm biến đổi HSV, tịnh tiến,
co giãn, lật ngang và mosaic.

Gợi ý phối đồ khác với tìm kiếm sản phẩm tương tự vì hệ thống phải đánh giá
khả năng bổ sung giữa các món đồ thuộc nhiều nhóm khác nhau~\cite{McAuley2015}. Han và cộng sự
mô hình hoá một bộ trang phục như một chuỗi và học đồng thời quan hệ tương
thích cùng biểu diễn thị giác--ngữ nghĩa~\cite{Han2017}. Trong mã nguồn hiện
tại, luồng gợi ý là một chức năng không dùng ảnh và không nhận kết quả YOLO.
Backend gửi hồ sơ cơ thể, giới tính tự khai, thời tiết hoặc mùa/nhiệt độ,
phong cách và hoàn cảnh đến Gemini; kết quả có một nhận định, lý do và từ một
đến ba điều chỉnh. Luồng hỏi đáp dùng cùng bối cảnh và tối đa sáu lượt gần
nhất khi gọi mô hình.

Thử trang phục ảo có thể được thực hiện bằng mô hình sinh ảnh hoặc bằng ghép
ảnh hai chiều. VITON là một phương pháp học sâu chuyển một món đồ lên ảnh
người theo chiến lược từ thô đến tinh mà không sử dụng mô hình 3D
~\cite{VITON2018}. Ứng dụng hiện dùng MediaPipe Pose Landmarker ở backend để
trả landmark và các điểm neo cơ thể qua WebSocket. Flutter làm mượt tư thế,
đặt ảnh trang phục nền trong suốt theo nhóm cơ thể và vẫn cho phép người dùng
dịch chuyển, co giãn, xoay thủ công. Đây là ghép lớp theo tư thế, không phải
mô hình sinh ảnh nên không mô phỏng chính xác độ rủ, nếp gấp hay ánh sáng.

Ba phiên bản thực nghiệm dùng cùng họ YOLO thay vì ba kiến trúc độc lập. V1
khởi tạo trực tiếp từ \code{yolo26n.pt}; V2 fine-tune checkpoint V1 trên dữ
liệu được bổ sung một phần từ DeepFashion; V3 sử dụng toàn bộ 52.712 ảnh
DeepFashion và lặp lại ảnh train cũ để giảm quên các lớp phụ kiện. V2 và V3
dùng AdamW, batch 16, kích thước ảnh 640 và weight decay 0,0005. Vì vậy, mục
tiêu so sánh hai kiến trúc học sâu khác nhau chưa được đáp ứng bởi các tệp
huấn luyện hiện có.

\section{Công cụ và thư viện}

Các thành phần thực tế được đối chiếu từ \code{.metadata},
\code{pubspec.lock}, môi trường Python của backend và cấu hình huấn luyện,
sau đó tổng hợp ở Bảng~\ref{tab:congcu}. Phiên bản dùng khi huấn luyện trên
Colab được phân biệt với phiên bản đang cài trong backend triển khai.

\begin{table}[H]
\centering
\caption{Công cụ và thư viện sử dụng trong đồ án.}
\label{tab:congcu}
\footnotesize
\setlength{\tabcolsep}{4pt}
\begin{tabularx}{\textwidth}{>{\bfseries\color{daunavy}\raggedright\arraybackslash}p{2.8cm} >{\color{daumaroon}\raggedright\arraybackslash}p{3.4cm} >{\raggedright\arraybackslash}X}
\toprule
\rowcolor{gray100}\bfseries\color{daunavy}Thành phần & \bfseries\color{daumaroon}Phiên bản & \bfseries\color{daunavy}Vai trò \\
\midrule
Flutter / Dart & 3.47.1 stable / 3.13.1 & Xây dựng giao diện Android, điều hướng bốn tab Phân tích, Gợi ý, Hỏi đáp và Cá nhân \\
Firebase & Core 4.14.0; Auth 6.6.1; Firestore 6.9.0 & Khởi tạo ứng dụng, xác thực email/mật khẩu, lưu hồ sơ tài khoản và hồ sơ cơ thể theo người dùng \\
FastAPI / Uvicorn & 0.141.1 / 0.52.4 & Cung cấp REST API và WebSocket cho YOLO, đánh giá/gợi ý, thời tiết và nhận diện tư thế~\cite{FastAPIFeatures} \\
Ultralytics YOLO & 8.4.135 (backend) & Nạp \code{backend/models/best.pt}, phát hiện 21 lớp gốc và chuẩn hoá về tám nhóm dùng trong ứng dụng \\
PyTorch & 2.13.0 (backend); 2.11.0+cu130 (Colab) & Chạy suy luận trên CPU/CUDA và huấn luyện các phiên bản YOLO V1--V3 \\
Gemini / Google Gen AI & google-genai 2.25.0; Gemini 3.5 Flash-Lite & Đánh giá bộ đồ, sinh gợi ý có cấu trúc và trả lời hỏi đáp; không tạo kết quả giả khi dịch vụ lỗi \\
MediaPipe & 1.0.1 & Pose Landmarker xử lý khung hình camera và trả landmark/điểm neo cho thử đồ qua WebSocket~\cite{MediaPipePose} \\
Google Colab & GPU Tesla T4 & Chuẩn bị dữ liệu và huấn luyện mô hình YOLO trên môi trường trực tuyến \\
OpenCV / NumPy / Pillow & 5.0.0.93 / 2.5.2 / 12.3.0 & Giải mã, kiểm tra ảnh, trích màu và chuẩn bị đầu vào phía backend \\
Camera và xử lý ảnh Flutter & camera 0.12.1; image\_picker 1.2.3; image 4.9.2 & Chụp/chọn ảnh, chuyển đổi frame camera và ghép lớp trang phục \\
Kết nối và vị trí & http 1.6.0; WebSocket 3.0.3; geolocator 14.0.3 & Gọi REST API, truyền frame pose thời gian thực và lấy vị trí cho dữ liệu Open-Meteo \\
\bottomrule
\end{tabularx}
\end{table}

Phiên bản Ultralytics 8.4.135 ở Bảng~\ref{tab:congcu} là phiên bản trong
môi trường backend cục bộ. Notebook Colab không in phiên bản Ultralytics nên
báo cáo không đồng nhất hai môi trường này một cách suy đoán.

\section{Các công trình liên quan}

YOLO đặt bài toán phát hiện đối tượng thành một quy trình hồi quy thống nhất,
cho phép dự đoán hộp giới hạn và lớp trong một lần chạy mạng
~\cite{Redmon2016}. Đồ án sử dụng một checkpoint YOLO kích thước nano để ưu
tiên khả năng triển khai, sau đó fine-tune trên dữ liệu thời trang riêng.

DeepFashion cung cấp dữ liệu thời trang có chú thích phong phú cho các bài
toán nhận diện, truy hồi và phân tích thuộc tính quần áo
~\cite{DeepFashion2016}. Công trình cho thấy dữ liệu thời trang cần chú ý
đến biến thiên về tư thế, góc chụp, phông nền và sự khác biệt nhỏ giữa các
nhóm sản phẩm.

Han và cộng sự nghiên cứu khả năng tương thích thời trang bằng Bi-LSTM,
kết hợp đặc trưng ảnh với thông tin ngữ nghĩa để hỗ trợ gợi ý món đồ và toàn
bộ trang phục~\cite{Han2017}. Hướng tiếp cận này có khả năng học quan hệ từ
dữ liệu, trong khi đồ án hiện mô tả phạm vi gợi ý dựa trên loại, màu sắc và
phong cách cơ bản.

VITON tạo ảnh thử đồ bằng một quy trình học sâu từ thô đến tinh
~\cite{VITON2018}. Đồ án không triển khai mô hình sinh ảnh tương đương mà
chọn kỹ thuật ghép lớp 2D đơn giản hơn để phù hợp thời gian và tài nguyên.
Sự khác biệt này cần được nêu rõ khi đánh giá: chức năng thử đồ của sản phẩm
chỉ mang tính minh hoạ trực quan, không phải mô phỏng chân thực.
